\documentclass[12pt,a4paper]{article}
\usepackage[utf8]{inputenc}
\usepackage[T1]{fontenc}
\usepackage[spanish]{babel}
\usepackage{amsmath}
\usepackage{amssymb}
\usepackage{geometry}
\usepackage{xcolor}
\usepackage{hyperref}
\usepackage{bookmark}
\usepackage{fancyhdr}
\usepackage{enumitem}
\usepackage{booktabs}
\usepackage{titlesec}
\geometry{margin=2.5cm}

% Encabezado estilo notas de clase
\setlength{\headheight}{15pt}
\pagestyle{fancy}
\fancyhf{}
% \fancyhead[L]{Fundamentos de Valoración de Activos}
% \fancyhead[R]{Exam}
\fancyfoot[C]{\thepage}
\fancyfoot[L]{TA Session}
\fancyfoot[R]{Jesus Villota Miranda}
\renewcommand{\footrulewidth}{0.4pt}

\newcommand{\R}{\mathbb{R}}

% Formato de ejercicios
\newenvironment{ejercicio}[1]{\subsection*{Ejercicio #1}}{}
\newenvironment{solucion}{\paragraph{Solución.}}{}

\begin{document}

\section{Rentas}

\begin{ejercicio}{1 (Renta semestral creciente diferida)}
Considere una renta semestral perpetua cuyo primer pago es $\alpha$ en el segundo año ($t=2$) y crece geométricamente a una tasa anual compuesta $\beta$, se descuenta con un tipo efectivo trimestral compuesto $\gamma$.
Asumiendo que $(1+\beta)^{1 / 2}<(1+\gamma)^2$, calcule el valor de la renta en $t=0$ en función de $\alpha$, $\beta$ y $\gamma$.

\medskip
\begin{enumerate}[label=\alph*)]
  \item $V=\dfrac{\alpha}{\left[(1+\gamma)^2-(1+\beta)^{1 / 2}\right]}$
  \item $\boxed{V=\dfrac{\alpha}{\left[(1+\gamma)^2-(1+\beta)^{1 / 2}\right] (1+\gamma)^6}}$
  \item $V=\dfrac{\alpha}{\left[(1+\gamma)^{1.5}-(1+\beta)^{1 / 2}\right] (1+\gamma)^{4.5}}$
\end{enumerate}
\end{ejercicio}

\begin{solucion}
\textbf{Notación y periodicidades.} \; Usaremos:
\begin{itemize}
  \item $A:=\alpha$ es el primer pago (en valor nominal).
  \item $g_1:=\beta$ es el crecimiento efectivo \emph{anual}. Necesitamos llevarlo a semestres.
  \item $i_4:=\gamma$ es el tipo efectivo \emph{trimestral}. Necesitamos llevarlo a semestres.
\end{itemize}

\textbf{Paso 1 (convertir crecimiento anual a crecimiento semestral).} \;
Los pagos son semestrales; queremos el crecimiento por semestre $g_2$ que, aplicado dos veces, reproduzca $g_1$:
\[
(1+g_2)^2=1+g_1
\quad\Rightarrow\quad
g_2=(1+g_1)^{1/2}-1.
\]

\textbf{Paso 2 (convertir descuento trimestral a descuento semestral).} \;
Un semestre contiene 2 trimestres; el tipo efectivo semestral $i_2$ equivalente a $i_4$ verifica:
\[
(1+i_2)=(1+i_4)^2
\quad\Rightarrow\quad
i_2=(1+i_4)^2-1.
\]

\textbf{Paso 3 (valorar la perpetuidad creciente en el instante adecuado).} \;
La renta es una \textbf{perpetuidad creciente} con pagos semestrales y primer pago en $t=2$.
El valor justo un semestre antes del primer pago ($t=1.5$) es:
\[
V_{1.5}=\frac{A}{i_2-g_2},
\qquad \text{(requiere } i_2>g_2\text{)}.
\]

\textbf{Paso 4 (descontar hasta $t=0$).} \;
Desde $t=0$ hasta $t=1.5$ hay $3$ semestres, así que:
\[
V_0=\frac{V_{1.5}}{(1+i_2)^3}
=\frac{A}{(i_2-g_2)(1+i_2)^3}.
\]
Sustituyendo $i_2=(1+i_4)^2-1$ y $g_2=(1+g_1)^{1/2}-1$:
\[
\boxed{
V_0=\frac{\alpha}{\left[(1+\gamma)^2-(1+\beta)^{1/2}\right]\,(1+\gamma)^6}
}
\]
donde hemos usado que $(1+i_2)^3=((1+i_4)^2)^3=(1+i_4)^6$.

\medskip
\textbf{Cómo leer las opciones:}
\begin{itemize}
  \item \textbf{a)} coincide con $V_{1.5}=\dfrac{A}{i_2-g_2}$: calcula el valor una semestralidad antes del primer pago, ya en periodicidad semestral (sin descontar a $t=0$).
  \item \textbf{b)} incluye el descuento adicional de 3 semestres, es decir, $V_0=\dfrac{A}{(i_2-g_2)(1+i_2)^3}$, que es la respuesta buscada.
  \item \textbf{c)} reproduce la misma estructura pero interpretando erróneamente $\gamma$ como \emph{tipo efectivo cuatrimestral} (i.e., 3 periodos/año) y sin alinear periodicidades; por eso entrega un valor con bases $(1+\gamma)^{1.5}$ y $(1+\gamma)^{4.5}$ que no corresponde al enunciado.
\end{itemize}
\end{solucion}

\end{document}


